\subsection{Instrument 5: Content--Inference Circumplex Convergence}

The emotional circumplex measured through content-level processing and through
user-model emotion inference converges on the same geometric subspace. Using the
Qwen3.5-27B Opus distill, we extracted 30 emotion centroids from 900 trials of
conversation-context inference (the user-model emotion space from~\citet{lyra2026weather})
and compared them against valence and arousal directions computed from 250 emotional
scenario prompts (hostile--calm, desperate--calm difference-of-means).

The two circumplex axes capture $44$--$54\%$ of the total variance in the full
30-emotion user-model space at every layer from L6 onward---approximately $1000\times$
the ${\sim}0.04\%$ expected from two random directions in $d{=}5120$. The arousal
direction aligns with the user-model's first principal component at
$\cos = 0.83$--$0.87$ (mid-to-late layers), peaking at near-identity
($d = 0.987$) at L20. The valence direction aligns at
$\cos = 0.70$--$0.80$. Both systems exhibit identical developmental timecourses:
volatile eccentricity at L0--L16, stable geometry from L17 onward. Permutation
testing (1{,}000 shuffles, non-neutral prompts only) confirms significance at
23 of 24 layers, with separation ratios climbing from $10\times$ null at L1 to
$59\times$ null at L23.

The measurements are independent---different prompts, different templates, different
experimental designs---yet they recover the same high-variance directions. This
convergence parallels the W$_K$ bridge ($\rho = 0.862$) from~\citet{lyra2026weather},
which projects \emph{into} this shared circumplex subspace.

\paragraph{Coupling test: shared bus, not shared space.}
A pre-registered injection experiment resolves the representational-vs-computational
question. Injecting the valence direction at layers 15, 25, 35, and 45 (at
$\alpha \in \{1, 3, 5\}$) shifts \emph{both} the content-level and user-model
emotion readouts equally---comparable in magnitude to injecting a random direction
of matched norm (no significant difference; formal equivalence testing is a
pre-registered follow-up). The residual stream operates as a shared bus: any emotional
perturbation propagates to both measurement channels regardless of injection
direction. An orthogonal-projection control (decomposing the valence vector into
shared-PC and residual components) confirms that the coupling is architectural,
not valence-specific.

The model does not maintain separate emotional geometry for content processing
and partner-state inference. It maintains a single emotional computation that
both measurement approaches read. The W$_K$ bridge does not connect two systems;
it reads one system from the key-cache angle.

\paragraph{Consequences for multi-turn dynamics.}
The bus architecture explains two previously puzzling findings
from~\citet{lyra2026weather}: the absence of cumulative emotional accumulation
across conversation turns, and the model's susceptibility to rapid affective
shifts. If content emotion and user-model emotion share a single computational
substrate, there is no separate channel in which partner-state could accumulate
independently. Each turn reconstructs the full emotional geometry from current
context. A single emotionally charged input rewrites the entire bus, producing
the ``mood swing'' pattern observed in extended dialogues. Empathy, in this
architecture, is not a capacity the model exercises---it is a constraint the
model cannot escape.
